\documentclass[11pt,a4paper]{article}
\usepackage[T1]{fontenc}
\usepackage{lmodern}
\usepackage[margin=27mm]{geometry}
\usepackage{amsmath,amssymb,amsthm,mathtools}
\usepackage{microtype}
\usepackage{booktabs,array}
\usepackage{enumitem}
\usepackage[hidelinks]{hyperref}
\setlength{\parskip}{5pt plus 1pt}
\setlength{\parindent}{0pt}
\setlength{\emergencystretch}{3em}
\setlist{itemsep=2pt,topsep=3pt}
\raggedbottom
\newtheorem{theorem}{Theorem}[section]
\newtheorem{lemma}[theorem]{Lemma}
\newtheorem{corollary}[theorem]{Corollary}
\newtheorem{proposition}[theorem]{Proposition}
\theoremstyle{definition}
\newtheorem{definition}[theorem]{Definition}
\newtheorem{example}[theorem]{Example}
\theoremstyle{remark}
\newtheorem{remark}[theorem]{Remark}
\newcommand{\Z}{\mathbb{Z}}
\newcommand{\N}{\mathbb{N}}
\newcommand{\Q}{\mathbb{Q}}
\newcommand{\F}{\mathbb{F}}
\newcommand{\ps}[1]{#1[[x]]}
\newcommand{\coeff}[2]{[x^{#1}]\,#2}
\newcommand{\hl}[1]{\textbf{#1}}
\newcommand{\para}[1]{\par\medskip\noindent\textbf{#1}\enspace}
\hypersetup{pdftitle={Odd coefficients, remainders modulo 9, and binary digits for the series of OEIS A396843},pdfauthor={Dev D. Goyal},pdfsubject={OEIS A396843}}
\title{Odd coefficients, remainders modulo $9$, and binary digits for the series $A(xA-3xA^2)=x^2$\\[6pt]\large A proof of one conjecture in OEIS A396843, and the correct form of another}
\author{Dev D. Goyal}
\date{October 2026}
\begin{document}
\maketitle
\begin{abstract}
Entry A396843 of the On-Line Encyclopedia of Integer Sequences (OEIS), contributed by Paul~D.~Hanna, lists the
coefficients $a(n)$ of the power series $A(x)=x+3x^2+15x^3+117x^4+\cdots$ determined by the equation
$A\bigl(xA(x)-3xA(x)^2\bigr)=x^2$. The entry conjectures that all coefficients are odd, and states a rule for their
remainders modulo $9$ that has since been found to be false. We prove that all coefficients are odd, and we determine
the remainders modulo $9$ exactly: for $n\ge2$, the remainder is $3$ if $n=4^k+1$, it is $6$ if $n=2\cdot4^k+1$, and it
is $0$ otherwise. We also show that $a(n)$ is divisible by $3^s$, where $s$ is the number of ones in the binary expansion
of $n-1$. All three results hold, in a form that depends on $m$, for the whole family of equations
$A\bigl(xA-mxA^2\bigr)=x^2$. The results were checked by computer and proved formally in the Lean proof assistant.
\end{abstract}
\begin{center}\begin{minipage}{0.88\linewidth}\small
\textbf{How this paper was produced.} This paper reports the result of an AI research effort directed by
Dev~D.~Goyal. ChatGPT (OpenAI), using the GPT-6 model with Pro reasoning and working autonomously from prompts by
Dev~D.~Goyal, found the proof, wrote its Lean~4 formalization, and wrote a first draft of the paper.
Dev~D.~Goyal chose the problems and asked for proofs, formal verification and revisions, but did not supply proof
ideas. Claude Opus~5.5, used with Claude Code, then checked the work independently: it rebuilt the formal proofs from
source and confirmed that they rely only on Lean's three standard axioms, recomputed the sequence from its definition
compared it with the published terms and tested every claim numerically, and read every step of the proofs. It also wrote the present
version of this paper.
Dev~D.~Goyal has verified the results and is responsible for them. The paper has not been peer reviewed.
The section ``About this paper'' at the end gives the details.
\end{minipage}\end{center}

\section{Introduction}

\subsection{The sequence}\label{sec:sequence}
The \emph{On-Line Encyclopedia of Integer Sequences} (OEIS)~\cite{oeisdb} is a large public database of integer
sequences. Many
entries also record \emph{conjectures}: patterns that the contributor observed in the data but did not prove.

This paper is about entry A396843~\cite{oeis3}, contributed by Paul~D.~Hanna in July 2026.
The entry defines a power series
\[
A(x)=a(1)\,x+a(2)\,x^2+a(3)\,x^3+\cdots=x+3x^2+15x^3+117x^4+975x^5+\cdots
\]
by a single equation,
\begin{equation}\label{eq:entry}
A\bigl(x\,A(x)-3\,x\,A(x)^2\bigr)=x^2 .
\end{equation}
To read the left side, first form the series $T(x)=x\,A(x)-3\,x\,A(x)^2$, which begins with $x^2$, and
then substitute $T(x)$ for $x$ in $A$. (Section~\ref{sec:background} explains what substituting one series into another
means.) The equation states that the result is $x^2$. The entry reads:
\begin{quote}\small\ttfamily\raggedright G.f. A(x) satisfies A( x*A(x) - 3*x*A(x)\^{}2 ) = x\^{}2.\par\smallskip Conjecture: all terms are odd.\par\smallskip Conjecture: for n > 3, a(n) == 1 (mod 9) iff n is of the form 4\^{}k+1, otherwise a(n) is divisible by 9. [This conjecture is false: n=5=4\^{}1+1, but a(5)=975 == 3 (mod 9), not 1. - Danial Asaria, Jul 26 2026]\par\smallskip\hfill\textrm{--- OEIS A396843, Paul D. Hanna, July 2026}\end{quote}

Here ``G.f.''\ stands for \emph{generating function}, the series whose coefficients are the terms of the sequence.

Equation~\eqref{eq:entry} determines the coefficients of $A$ one at a time, as follows. Write $A=x+a(2)x^2+\cdots$, with first coefficient $a(1)=1$ as in the entry. Then
\[
T=x\,A-3\,x\,A^2=x^2+\bigl(a(2)-3\bigr)x^3+\cdots,
\qquad
A(T)=T+a(2)\,T^2+\cdots=x^2+\bigl(a(2)-3\bigr)x^3+\cdots,
\]
because $T^2$ begins with $x^4$. For this to equal $x^2$, the coefficient of $x^3$ must vanish, so
$a(2)=3$. In the same way, the coefficient of $x^{n+1}$ in \eqref{eq:entry} determines $a(n)$ once
$a(1),\dots,a(n-1)$ are known (Lemma~\ref{lem:exist}). So there is exactly one solution, and all its coefficients are
integers. Its first terms are 1, 3, 15, 117, 975, 8631, 79839, 763209, 7477701, \dots{} (starting from $n=1$).

\subsection{The question}
The entry contains two conjectures (see the quotation above). The first says that every $a(n)$ is odd. The
second describes the remainders of $a(n)$ on division by $9$. (Section~\ref{sec:background} recalls the notation
$a\equiv b\pmod m$ for ``$a$ and $b$ leave the same remainder on division by $m$''.) As the note by Danial~Asaria in the
entry records, the second conjecture is false as stated: $5=4^1+1$, but $a(5)=975$ leaves remainder $3$, not $1$, on
division by $9$. So the questions are: are all terms odd, and what is the correct rule modulo $9$?

The table shows the remainders of the first terms on division by $2$ and by $9$. The indices in bold are those of the
form $n=2^j+1$, namely $2,3,5,9,17,33$.
\begin{center}\small\setlength{\tabcolsep}{3.2pt}
\begin{tabular}{>{\raggedright\arraybackslash}p{6.4em}>{\centering\arraybackslash}p{1.15em}>{\centering\arraybackslash}p{1.15em}>{\centering\arraybackslash}p{1.15em}>{\centering\arraybackslash}p{1.15em}>{\centering\arraybackslash}p{1.15em}>{\centering\arraybackslash}p{1.15em}>{\centering\arraybackslash}p{1.15em}>{\centering\arraybackslash}p{1.15em}>{\centering\arraybackslash}p{1.15em}>{\centering\arraybackslash}p{1.15em}>{\centering\arraybackslash}p{1.15em}>{\centering\arraybackslash}p{1.15em}>{\centering\arraybackslash}p{1.15em}>{\centering\arraybackslash}p{1.15em}>{\centering\arraybackslash}p{1.15em}>{\centering\arraybackslash}p{1.15em}>{\centering\arraybackslash}p{1.15em}}
\toprule
$n$ & 1 & \hl{2} & \hl{3} & 4 & \hl{5} & 6 & 7 & 8 & \hl{9} & 10 & 11 & 12 & 13 & 14 & 15 & 16 & \hl{17} \\ \midrule
$a(n)\bmod 2$ & 1 & \hl{1} & \hl{1} & 1 & \hl{1} & 1 & 1 & 1 & \hl{1} & 1 & 1 & 1 & 1 & 1 & 1 & 1 & \hl{1} \\
$a(n)\bmod 9$ & 1 & \hl{3} & \hl{6} & 0 & \hl{3} & 0 & 0 & 0 & \hl{6} & 0 & 0 & 0 & 0 & 0 & 0 & 0 & \hl{3} \\
\bottomrule
\end{tabular}
\\[4pt]
\begin{tabular}{>{\raggedright\arraybackslash}p{6.4em}>{\centering\arraybackslash}p{1.15em}>{\centering\arraybackslash}p{1.15em}>{\centering\arraybackslash}p{1.15em}>{\centering\arraybackslash}p{1.15em}>{\centering\arraybackslash}p{1.15em}>{\centering\arraybackslash}p{1.15em}>{\centering\arraybackslash}p{1.15em}>{\centering\arraybackslash}p{1.15em}>{\centering\arraybackslash}p{1.15em}>{\centering\arraybackslash}p{1.15em}>{\centering\arraybackslash}p{1.15em}>{\centering\arraybackslash}p{1.15em}>{\centering\arraybackslash}p{1.15em}>{\centering\arraybackslash}p{1.15em}>{\centering\arraybackslash}p{1.15em}>{\centering\arraybackslash}p{1.15em}>{\centering\arraybackslash}p{1.15em}}
\toprule
$n$ & 18 & 19 & 20 & 21 & 22 & 23 & 24 & 25 & 26 & 27 & 28 & 29 & 30 & 31 & 32 & \hl{33} & 34 \\ \midrule
$a(n)\bmod 2$ & 1 & 1 & 1 & 1 & 1 & 1 & 1 & 1 & 1 & 1 & 1 & 1 & 1 & 1 & 1 & \hl{1} & 1 \\
$a(n)\bmod 9$ & 0 & 0 & 0 & 0 & 0 & 0 & 0 & 0 & 0 & 0 & 0 & 0 & 0 & 0 & 0 & \hl{6} & 0 \\
\bottomrule
\end{tabular}
\end{center}
All terms in the table are odd, and from $n=2$ on, the remainder modulo $9$ is $0$ except at the indices in
bold, where it is $3$ or $6$.

\subsection{What this paper proves}
\textbf{Main results} (Theorems~\ref{thm:parity}, \ref{thm:modsq}, \ref{thm:binary} and Corollary~\ref{cor:entry}).
\begin{enumerate}
\item \emph{Every $a(n)$ is odd.} This proves the first conjecture.
\item \emph{For $n\ge2$, the remainder of $a(n)$ on division by $9$ is $3$ if $n=4^k+1$ (that is, $n=2,5,17,65,\dots$), it is
$6$ if $n=2\cdot4^k+1$ ($n=3,9,33,129,\dots$), and it is $0$ for all other $n$.} This is the correct form of the second
conjecture. It shows that the conjecture fails in two ways: at $n=4^k+1$ with $k\ge1$ the remainder is $3$, not $1$; and at
$n=2\cdot4^k+1$ with $k\ge1$ the number $a(n)$ is not divisible by $9$.
\item \emph{$a(n)$ is divisible by $3^{s_2(n-1)}$, where $s_2(n-1)$ is the number of ones in the binary expansion of
$n-1$.} For example, $3^4=81$ divides $a(16)$, because $15$ is $1111$ in binary. The same holds for the coefficients of
the compositional inverse of $A$.
\end{enumerate}

All of these statements are special cases of theorems about the whole family of equations
\[
A_m\bigl(x\,A_m(x)-m\,x\,A_m(x)^2\bigr)=x^2,
\]
in which $3$ is replaced by an arbitrary integer $m$. Hanna's entries A265940, A271958, A396843,
A396844 and A396845~\cite{oeis1,oeis2,oeis3,oeis4,oeis5} are the cases $m=1,2,3,4,5$. The central result is a
formula for the whole series modulo $m^2$ (Theorem~\ref{thm:modsq}):
\[
A_m(x)\equiv x+m\,\bigl(x^2-x^3+x^5-x^9+x^{17}-x^{33}+\cdots\bigr)\pmod{m^2},
\]
where the exponents are the numbers $2^j+1$ and the signs alternate.

\subsection{What was known}
The entry records several formulas due to Hanna. One of them, $B(x^2)=x\,A(x)-3\,x\,A(x)^2$, where $B$ is
the compositional inverse of $A$ (Section~\ref{sec:background}), is the starting point of our proof of the divisibility
result. In July 2026, Danial~Asaria added to the entry the counterexample $a(5)=975$ to the conjecture modulo~$9$.
We are not aware of an earlier proof of these statements.

\subsection{Motivation}
Equation~\eqref{eq:entry} mixes two operations, multiplication and substitution. Such equations rarely have explicit
solutions: the coefficients can be computed one after another, but there is usually no formula for them, and they grow
quickly. Hanna has contributed many sequences of this kind to the OEIS. For such sequences, divisibility patterns are
often the only structure that is visible in the data, and proving them is one way to understand the sequence.

In the divisibility result, the binary digits of $n-1$ control the power of $m$ that divides $a(n)$. Statements of
this kind are classical for explicit numbers. For instance, Legendre's formula says that the exponent of $2$
in $n!$ equals $n-s_2(n)$, where $s_2(n)$ is the number of ones in the binary expansion of $n$~\cite{concrete}.
Here the same kind of structure appears in a series that is defined only implicitly.

Finally, the proofs use a few short ideas that recur throughout this collection of papers: reduce the equation modulo
a small number, guess the solution, and conclude by uniqueness; and use the compositional inverse to turn an implicit
equation into a recursion.

\subsection{How to read this paper}
Section~\ref{sec:background} explains the notions we need: power series, substitution, compositional inverses,
congruences, binary digits and prime exponents. Section~\ref{sec:results} states the results precisely, and
Section~\ref{sec:proofs} proves them. Section~\ref{sec:checks} describes how the results were checked, by computer and by
a formal proof in Lean. Section~\ref{sec:discussion} discusses the scope of the results and open questions.
\section{Background}\label{sec:background}
This section collects the definitions and facts used later. Readers who know formal power series and congruences can skip it and return when needed. Standard references for formal power series are~\cite{niven,stanley}.
\subsection{Formal power series}
A \emph{formal power series} is an expression
\[
U(x)=u_0+u_1x+u_2x^2+u_3x^3+\cdots ,
\]
determined by its sequence of coefficients $u_0,u_1,u_2,\dots$; see~\cite{niven,stanley}. No number is ever
substituted for $x$, so questions of convergence do not arise. We write $\coeff{n}{U}=u_n$ for the coefficient of $x^n$ in $U$; two series are equal
when all their coefficients are equal.

Series are added coefficient by coefficient and multiplied like polynomials:
\[
\coeff{n}{(UV)}=u_0v_n+u_1v_{n-1}+\cdots+u_nv_0 .
\]
Each coefficient of a product is a finite sum. For example, $(1-x)(1+x+x^2+x^3+\cdots)=1$, so $1+x+x^2+\cdots$ is
the \emph{multiplicative inverse} of $1-x$, written $\frac{1}{1-x}$. More generally, every series $D$ with constant term
$1$ has a unique multiplicative inverse $\frac1D$, whose coefficients are determined one at a time by $D\cdot\frac1D=1$
using only addition, subtraction and multiplication. Whenever a quotient of power series appears in this paper, its
denominator has constant term $1$, so the quotient is again a power series.

A \emph{commutative ring} is a set with an addition, a subtraction and a multiplication that satisfy the usual rules
of arithmetic; examples are $\Z$, $\Q$ and the integers modulo $m$. For a commutative ring $R$ we write $\ps{R}$ for
the set of power series with coefficients in $R$; for instance, writing $\F_2$ for the integers modulo $2$, $\ps{\F_2}$
consists of the series with coefficients modulo $2$. We write
$U(0)$ for the constant term $u_0$ of $U$, and $U=O(x^d)$ when the coefficients of $1,x,\dots,x^{d-1}$ in $U$ are all
zero. Thus $U-V=O(x^d)$ means that $U$ and $V$ have the same coefficient of $x^n$ for every $n<d$. A series with
constant term $0$ and coefficient of $x$ equal to $1$, that is, a series of the form $x+u_2x^2+u_3x^3+\cdots$, is
called \emph{normalized}.

\subsection{Substituting one series into another}
If $V(x)$ has constant term $0$, we can substitute $V$ for $x$ in $U$:
\[
U(V(x))=u_0+u_1V(x)+u_2V(x)^2+u_3V(x)^3+\cdots .
\]
This is called \emph{composition}, also written $U\circ V$. It makes sense because $V(x)^k$ has no terms below $x^k$,
so only the finitely many terms with $k\le n$ affect the coefficient of $x^n$. For example, if $U=x+x^2$ and $V=x^2$, then
$U(V(x))=x^2+x^4$. Composition is associative: $U\circ(V\circ W)=(U\circ V)\circ W$.

Substituting a series into itself repeatedly gives its \emph{iterates}: $U^{\circ2}(x)=U(U(x))$,
$U^{\circ3}(x)=U(U(U(x)))$, and so on, with $U^{\circ0}(x)=x$ and $U^{\circ1}=U$. They must not be confused with
\emph{powers} $U^2=U\cdot U$, $U^3=U\cdot U\cdot U$. (The OEIS sometimes writes $A^k(x)$ for the $k$-th iterate. In
this paper the small circle always means an iterate, and a plain exponent always means a power.)

\subsection{Undoing a substitution}
A normalized series can always be ``undone'' by another normalized series.

\begin{lemma}[Compositional inverse]\label{lem:inverse}
For every normalized series $U$ there is exactly one normalized series $V$ with $U(V(x))=x$, and this $V$ also
satisfies $V(U(x))=x$. If $U$ has integer coefficients, so does $V$.
\end{lemma}

We call $V$ the \emph{compositional inverse} of $U$. The OEIS calls it the \emph{series reversion} of $U$. For example,
the compositional inverse of $x-x^2$ is $C(x)=x+x^2+2x^3+5x^4+14x^5+\cdots$, whose coefficients are the Catalan numbers.
Indeed $C-C^2=x$, which says exactly that substituting $C$ into $x-x^2$ gives $x$.

\begin{proof}
Write $V=x+v_2x^2+v_3x^3+\cdots$ with unknown coefficients. In
$U(V)=V+u_2V^2+u_3V^3+\cdots$, the coefficient of $x^n$ (for $n\ge2$) equals $v_n$ plus terms that involve only
$v_2,\dots,v_{n-1}$. Setting it equal to $0$ therefore determines $v_n$ from the earlier coefficients, using only
additions, subtractions and multiplications. So there is exactly one $V$ with $U\circ V=x$, and it has integer
coefficients if $U$ does. Applying the same construction to $V$ gives a normalized $W$ with $V\circ W=x$. Then
$U=U\circ(V\circ W)=(U\circ V)\circ W=W$, so $V\circ U=V\circ W=x$ as well.
\end{proof}

\subsection{Remainders and congruences}
For an integer $m\ge2$ we write $a\equiv b\pmod m$ (``$a$ is congruent to $b$ modulo $m$'') when $a-b$ is divisible by
$m$, that is, when $a$ and $b$ leave the same remainder after division by $m$. The remainder of $a$ on division by
$m$ is called its \emph{residue} modulo $m$; it is one of $0,1,\dots,m-1$.

Addition and multiplication are compatible with congruences, so the residues $0,1,\dots,m-1$, added and multiplied
modulo $m$, form a ring, written $\Z/m\Z$. When $m=p$ is a prime we also write $\F_p$. In $\F_2=\{0,1\}$ we have $1+1=0$,
so addition and subtraction coincide.

For power series with integer coefficients, $U\equiv V\pmod m$ means that \emph{every} coefficient of $U-V$ is
divisible by $m$. Equivalently, the two series become equal when every coefficient is replaced by its residue; we call
this \emph{reducing} the series modulo $m$. The key fact is the following.

\emph{Reduction modulo $m$ is compatible with all the operations above.} To reduce a sum, a product, a substitution, or
an inverse, we may reduce the ingredients first and then compute in $\Z/m\Z$. This holds because every
coefficient of the result is obtained from finitely many coefficients of the ingredients by adding, subtracting and
multiplying, and these operations are compatible with remainders.

\para{Rings.} Some statements below hold over every commutative ring, for example $\Z$, $\Q$ and $\Z/m\Z$. No
division is used in them; this matters because in $\Z/m\Z$ division is not always possible (for instance, $2$ has no inverse in
$\Z/4\Z$).

\subsection{The method}
The main proofs follow three steps.
\begin{enumerate}
\item \emph{Uniqueness.} Show that the defining equation determines the coefficients one at a time, using only
addition, subtraction and multiplication. Then the equation has exactly one solution of the required shape, in every
commutative ring, including $\Z/m\Z$.
\item \emph{Reduction.} Reducing the integer solution modulo $m$ gives a solution of the reduced equation,
because reduction is compatible with all operations.
\item \emph{A candidate solution.} Find an explicit series that satisfies the reduced equation. By uniqueness, it is
the reduction of the integer solution, and its coefficients can be read off.
\end{enumerate}
Finding the candidate in step~3 is usually the main difficulty.
\subsection{The Frobenius identity}
\begin{lemma}[Frobenius identity]\label{lem:frobenius}
Let $p$ be a prime. For every power series $U$ with coefficients in $\F_p$, we have $U(x)^p=U(x^p)$.
\end{lemma}

\begin{proof}
In $\F_p$, the binomial theorem simplifies to $(a+b)^p=a^p+b^p$, because every binomial coefficient
$\binom pi$ with $0<i<p$ is divisible by $p$. Also $c^p=c$ for every $c\in\F_p$ (Fermat's little theorem). So
$\bigl(\sum_nu_nx^n\bigr)^p=\sum_nu_n^px^{pn}=\sum_nu_nx^{pn}=U(x^p)$. The first equality holds for polynomials by
repeated use of $(a+b)^p=a^p+b^p$; it holds for power series because each coefficient of $U(x)^p$ involves only finitely
many coefficients of $U$.
\end{proof}

For $p=2$ this says that squaring a series modulo $2$ doubles every exponent: for instance,
$(x+x^3)^2\equiv x^2+x^6\pmod2$. In particular, the square of any series has, modulo $2$, no terms of odd degree.

\subsection{Binary digits}
Every integer $r\ge0$ can be written in exactly one way as a sum of distinct powers of $2$; this is its binary
expansion. For example $13=8+4+1$, which is written $1101$ in binary. We write $s_2(r)$ for the number of powers of $2$
in this sum, that is, the number of ones in the binary expansion of $r$. So $s_2(0)=0$, $s_2(13)=3$, $s_2(2^k)=1$, and
$s_2(2^k-1)=k$, because $2^k-1=1+2+4+\cdots+2^{k-1}$. We use two facts:
\begin{equation}\label{eq:digits}
s_2(2r)=s_2(r),\qquad s_2(r_1+r_2+\cdots+r_k)\le s_2(r_1)+s_2(r_2)+\cdots+s_2(r_k).
\end{equation}
The first holds because doubling a number shifts its binary digits one place to the left. For the second, write each
$r_i$ as a sum of distinct powers of $2$ and put all these sums together. This writes the total as a sum of
$s_2(r_1)+\cdots+s_2(r_k)$ powers of $2$, possibly with repetitions. Now repeatedly replace two equal powers
$2^i+2^i$ by the single power $2^{i+1}$. This does not change the value, and it lowers the number of terms by one.
When no repetitions are left, we have the binary expansion of the total, with $s_2(r_1+\cdots+r_k)$ terms.

\subsection{Prime exponents}
For a prime number $p$ and a nonzero integer $z$, we write $v_p(z)$ for the exponent of $p$ in the prime factorization
of $z$, that is, the largest $e$ such that $p^e$ divides $z$. For example $975=3\cdot5^2\cdot13$, so $v_3(975)=1$ and
$v_5(975)=2$. Every power of $p$ divides $0$, so we set $v_p(0)=\infty$.

\section{Results}\label{sec:results}

We treat the whole family at once. For an integer $m$, consider the equation
\begin{equation}\label{eq:general}
A_m\bigl(x\,A_m(x)-m\,x\,A_m(x)^2\bigr)=x^2,\qquad A_m(x)=x+(\text{terms of degree}\ge2).
\end{equation}
Lemma~\ref{lem:exist} in Section~\ref{sec:proofs} shows that for every integer $m$ it has exactly one solution, and
that the solution has integer coefficients. We write $a_n(m)$ for the coefficient of $x^n$ in $A_m$. The sequence of
OEIS A396843 is $a(n)=a_n(3)$.

\begin{theorem}[Parity]\label{thm:parity}
If $m$ is odd, then $a_n(m)$ is odd for every $n\ge1$. If $m$ is even, then $a_1(m)=1$ and $a_n(m)$ is even for every
$n\ge2$.
\end{theorem}

\begin{theorem}[Remainders modulo $m^2$]\label{thm:modsq}
Let $m\ge2$. Then
\[
A_m(x)\equiv x+m\sum_{j\ge0}(-1)^j\,x^{2^j+1}=x+m\,\bigl(x^2-x^3+x^5-x^9+x^{17}-\cdots\bigr)\pmod{m^2}.
\]
In other words, for $n\ge2$: if $n=2^j+1$ for some $j\ge0$, then $a_n(m)\equiv(-1)^j\,m\pmod{m^2}$; for all other $n$,
$a_n(m)$ is divisible by $m^2$.
\end{theorem}

Let $B_m$ be the compositional inverse of $A_m$ (Lemma~\ref{lem:inverse}), so that $A_m(B_m(x))=x$, and let $b_n(m)$ be
its coefficient of $x^n$. For example, the entry records $B_{3}(x)=x-3x^2+3x^3-27x^4+105x^5-\cdots.$
Recall that $s_2(r)$ is the number of ones in the binary expansion of $r$.

\begin{theorem}[Binary digits control divisibility]\label{thm:binary}
Let $m\ge2$ and $n\ge1$.
\begin{enumerate}[label=(\alph*)]
\item $m^{\,s_2(n-1)}$ divides both $a_n(m)$ and $b_n(m)$. In particular, $m^k$ divides $a_{2^k}(m)$ for every $k\ge0$.
\item The same holds with $m$ replaced by a variable $t$. More precisely, there are polynomials $a_n(t)$ and $b_n(t)$
in $t$ with integer coefficients whose values at $t=m$ are $a_n(m)$ and $b_n(m)$ for every integer $m$; they are
constructed in Section~\ref{sec:proofs}. These polynomials have no terms of degree less than $s_2(n-1)$ in $t$.
\end{enumerate}
\end{theorem}

For example, $a_4(t)=5t^3-2t^2$, and $s_2(3)=2$. The next table illustrates part~(a) for $m=3$. It compares
$s_2(n-1)$ with the exponent $v_{3}(a(n))$ of the prime $3$ in $a(n)$. By
Corollary~\ref{cor:val} below, the second is always at least $s_2(n-1)$.
Indices in bold are the powers of $2$, $n=2^k$, where $s_2(n-1)=k$.
\begin{center}\small\setlength{\tabcolsep}{3.2pt}
\begin{tabular}{>{\raggedright\arraybackslash}p{6.4em}>{\centering\arraybackslash}p{1.15em}>{\centering\arraybackslash}p{1.15em}>{\centering\arraybackslash}p{1.15em}>{\centering\arraybackslash}p{1.15em}>{\centering\arraybackslash}p{1.15em}>{\centering\arraybackslash}p{1.15em}>{\centering\arraybackslash}p{1.15em}>{\centering\arraybackslash}p{1.15em}>{\centering\arraybackslash}p{1.15em}>{\centering\arraybackslash}p{1.15em}>{\centering\arraybackslash}p{1.15em}>{\centering\arraybackslash}p{1.15em}>{\centering\arraybackslash}p{1.15em}>{\centering\arraybackslash}p{1.15em}>{\centering\arraybackslash}p{1.15em}>{\centering\arraybackslash}p{1.15em}}
\toprule
$n$ & \hl{1} & \hl{2} & 3 & \hl{4} & 5 & 6 & 7 & \hl{8} & 9 & 10 & 11 & 12 & 13 & 14 & 15 & \hl{16} \\ \midrule
$s_2(n-1)$ & \hl{0} & \hl{1} & 1 & \hl{2} & 1 & 2 & 2 & \hl{3} & 1 & 2 & 2 & 3 & 2 & 3 & 3 & \hl{4} \\
$v_{3}(a(n))$ & \hl{0} & \hl{1} & 1 & \hl{2} & 1 & 2 & 3 & \hl{3} & 1 & 2 & 2 & 4 & 2 & 4 & 3 & \hl{5} \\
\bottomrule
\end{tabular}
\\[4pt]
\begin{tabular}{>{\raggedright\arraybackslash}p{6.4em}>{\centering\arraybackslash}p{1.15em}>{\centering\arraybackslash}p{1.15em}>{\centering\arraybackslash}p{1.15em}>{\centering\arraybackslash}p{1.15em}>{\centering\arraybackslash}p{1.15em}>{\centering\arraybackslash}p{1.15em}>{\centering\arraybackslash}p{1.15em}>{\centering\arraybackslash}p{1.15em}>{\centering\arraybackslash}p{1.15em}>{\centering\arraybackslash}p{1.15em}>{\centering\arraybackslash}p{1.15em}>{\centering\arraybackslash}p{1.15em}>{\centering\arraybackslash}p{1.15em}>{\centering\arraybackslash}p{1.15em}>{\centering\arraybackslash}p{1.15em}>{\centering\arraybackslash}p{1.15em}}
\toprule
$n$ & 17 & 18 & 19 & 20 & 21 & 22 & 23 & 24 & 25 & 26 & 27 & 28 & 29 & 30 & 31 & \hl{32} \\ \midrule
$s_2(n-1)$ & 1 & 2 & 2 & 3 & 2 & 3 & 3 & 4 & 2 & 3 & 3 & 4 & 3 & 4 & 4 & \hl{5} \\
$v_{3}(a(n))$ & 1 & 2 & 4 & 3 & 2 & 3 & 3 & 5 & 3 & 3 & 4 & 4 & 4 & 5 & 6 & \hl{6} \\
\bottomrule
\end{tabular}
\end{center}

\begin{corollary}[Exponents of the primes dividing $m$]\label{cor:val}
Let $m\ge2$, let $p$ be a prime that divides $m$, and let $e=v_p(m)$. Then $v_p(a_n(m))\ge e\,s_2(n-1)$ for every
$n\ge1$. Moreover, for $n\ge2$,
\[
v_p\bigl(a_n(m)\bigr)=e\quad\Longleftrightarrow\quad n=2^j+1\ \text{for some }j\ge0,
\]
and $v_p(a_n(m))>e$ for all other $n\ge2$.
\end{corollary}

\begin{corollary}[OEIS A396843]\label{cor:entry}
For $m=3$, that is, for the sequence $a(n)$ of OEIS A396843:
\begin{enumerate}[label=(\alph*)]
\item every $a(n)$ with $n\ge1$ is odd;
\item $3^{s_2(n-1)}$ divides $a(n)$ and $b_n(3)$, for every $n\ge1$;
\item for $n\ge2$,
\[
a(n)\equiv\begin{cases}3\pmod9,&\text{if }n=4^k+1\text{ for some }k\ge0,\\
6\pmod9,&\text{if }n=2\cdot4^k+1\text{ for some }k\ge0,\\
0\pmod9,&\text{otherwise.}\end{cases}
\]
\end{enumerate}
In particular, the second conjecture of the entry fails in two ways. At $n=4^k+1$ with $k\ge1$ the remainder is $3$,
not $1$; Asaria's counterexample $a(5)=975$ is the case $k=1$. At $n=2\cdot4^k+1$ with $k\ge1$ the remainder is $6$, so
$9$ does not divide $a(n)$; the first case is $a(9)=7477701$.
\end{corollary}

\section{Proofs}\label{sec:proofs}

\para{The idea in brief.} Each theorem has its own short idea.
\emph{Parity:} modulo $2$, the equation is solved by $x/(1+x)=x+x^2+x^3+\cdots$ when $m$ is odd, and by $x$ when $m$ is
even; since the solution is unique, this is the reduction of $A_m$.
\emph{Remainders modulo $m^2$:} modulo $m$ the solution is $x$, so $A_m=x+m\,C$ for some integer series $C$.
Putting this into the equation and discarding multiples of $m^2$ leaves the linear equation
$x\,C(x)+C(x^2)\equiv x^3\pmod m$, whose solution has nonzero terms only in the degrees $2^j+1$.
\emph{Binary digits:} substituting the inverse series into the equation gives an identity that expresses each
coefficient through earlier ones. In that identity, an index is either halved, which does not change its binary digit
sum, or split into two parts at the cost of one extra factor $m$. An induction on $n$ then shows the divisibility.

\begin{lemma}[Existence and uniqueness]\label{lem:exist}
Let $R$ be a commutative ring and let $m\in R$. There is exactly one series $A=x+a_2x^2+a_3x^3+\cdots$ with coefficients
in $R$ such that $A\bigl(xA-mxA^2\bigr)=x^2$. For $R=\Z$ this is the series $A_m$, and for every $q\ge2$ the reduction of
$A_m$ modulo $q$ is the solution over $\Z/q\Z$.
\end{lemma}

\begin{proof}
For a series $A=x+a_2x^2+\cdots$ write $T_A=xA-mxA^2$, which begins with $x^2$, and $\Psi(A)=A(T_A)$. The key
claim is: \emph{if $A$ and $A^*$ agree below degree $r$, where $r\ge2$, and $A^*-A=c\,x^r+(\text{higher terms})$, then
\begin{equation}\label{eq:perturb}
\Psi(A^*)-\Psi(A)=c\,x^{r+1}+(\text{higher terms}).
\end{equation}}
In other words, changing the coefficient of $x^r$ by $c$ changes the coefficient of $x^{r+1}$ in $\Psi$ by exactly $c$, and
leaves all lower coefficients alone. To see this, split the difference into two parts,
\[
\Psi(A^*)-\Psi(A)=\bigl(A(T_{A^*})-A(T_A)\bigr)+(A^*-A)(T_{A^*}).
\]
First, $T_{A^*}-T_A=x(A^*-A)\bigl(1-m(A+A^*)\bigr)=c\,x^{r+1}+(\text{higher terms})$. Next,
$A(T_{A^*})-A(T_A)=(T_{A^*}-T_A)+\sum_{k\ge2}a_k\,(T_{A^*}^k-T_A^k)$, and for $k\ge2$ the difference
$T_{A^*}^k-T_A^k$ is $T_{A^*}-T_A$ times a sum of products of $k-1$ factors, each beginning with $x^2$; so it begins
in degree $r+3$ or later. Hence $A(T_{A^*})-A(T_A)=c\,x^{r+1}+(\text{higher terms})$. Finally,
$(A^*-A)(T_{A^*})$ begins in degree $2r\ge r+2$, because $A^*-A$ begins in degree $r$ and $T_{A^*}$ in degree $2$.
This proves \eqref{eq:perturb}.

\emph{Uniqueness.} If two different solutions existed, let $r$ be the first degree in which they differ. By
\eqref{eq:perturb}, $\Psi$ of the two would differ in degree $r+1$. But both equal $x^2$.

\emph{Existence.} Start with $A=x$. If $\Psi(A)-x^2$ has no terms below degree $r+1$ and its coefficient of $x^{r+1}$
is $c$, replace $A$ by $A-c\,x^r$. By \eqref{eq:perturb} this removes that coefficient and changes nothing below it.
Each step fixes one more coefficient of $A$ for good, because the coefficient of $x^d$ in $\Psi(A)$ depends only on
$a_2,\dots,a_{d-1}$. So the steps build a solution, one coefficient at a time.

No division was used, so the argument works over every commutative ring. Reducing the integer solution modulo $q$
gives a solution over $\Z/q\Z$, because reduction is compatible with all the operations; by uniqueness it is
\emph{the} solution there.
\end{proof}

Running the construction with a variable $t$ in place of $m$, that is, over the ring $\Z[t]$ of polynomials in $t$ with
integer coefficients, gives
\[
A=x+t\,x^2+(2t^2-t)\,x^3+(5t^3-2t^2)\,x^4+\cdots,
\]
and in general a polynomial $a_n(t)$ with integer coefficients as the coefficient of $x^n$. Substituting the number $m$
for $t$ is compatible with all the operations, so by uniqueness it turns this series into $A_m$. So $a_n(m)$ is the value
of $a_n(t)$ at $t=m$; for example $a_3(3)=2\cdot9-3=15$. The same applies to the inverse series and $b_n(t)$.

\begin{proof}[Proof of Theorem~\ref{thm:parity}]
Let $m$ be odd and reduce modulo $2$. Since $m\equiv1$ and $-1\equiv1$, the equation becomes $P(xP+xP^2)=x^2$ over
$\F_2$. Consider $H=x/(1+x)=x+x^2+x^3+\cdots$. Working in $\ps{\F_2}$, where $(1+x)^2=1+x^2$ (Lemma~\ref{lem:frobenius}),
\begin{align*}
xH+xH^2&=\frac{x^2(1+x)+x^3}{(1+x)^2}=\frac{x^2+2x^3}{(1+x)^2}=\frac{x^2}{(1+x)^2},\\
H\!\left(\frac{x^2}{(1+x)^2}\right)&=\frac{x^2}{(1+x)^2+x^2}=\frac{x^2}{1+2x+2x^2}=x^2 ,
\end{align*}
because $2=0$ in $\F_2$. (For the second line we used $H(y)=y/(1+y)$ and multiplied the numerator and the denominator
by $(1+x)^2$.) So $H$ solves the reduced equation. By Lemma~\ref{lem:exist}, $A_m\equiv H\pmod2$: all
coefficients of positive degree are odd.

If $m$ is even, the reduced equation is $P(xP)=x^2$. It is solved by $P=x$, because $x\cdot x=x^2$. So $A_m\equiv
x\pmod2$, which is the claim.
\end{proof}

\begin{proof}[Proof of Theorem~\ref{thm:modsq}]
\emph{Step 1: modulo $m$.} Modulo $m$ the equation becomes $P(xP)=x^2$, which is solved by $P=x$. By
Lemma~\ref{lem:exist}, $A_m\equiv x\pmod m$. So every coefficient of $A_m-x$ is divisible by $m$, and we can write
$A_m=x+m\,C$ for a series $C$ with integer coefficients that begins in degree $2$.

\emph{Step 2: modulo $m^2$.} Substitute $A_m=x+mC$ into $T=xA_m-mxA_m^2$:
\begin{align*}
T&=x^2+m\,xC-m\,x\,(x+mC)^2=x^2+m\,(xC-x^3)-2m^2x^2C-m^3xC^2\\
&\equiv x^2+m\,(xC-x^3)\pmod{m^2}.
\end{align*}
Since $T\equiv x^2\pmod m$, also $C(T)\equiv C(x^2)\pmod m$, and so $m\,C(T)\equiv m\,C(x^2)\pmod{m^2}$. Therefore
\[
x^2=A_m(T)=T+m\,C(T)\equiv x^2+m\,\bigl(xC(x)-x^3+C(x^2)\bigr)\pmod{m^2}.
\]
So every coefficient of $m\bigl(xC(x)+C(x^2)-x^3\bigr)$ is divisible by $m^2$. Dividing by $m$, this says
\begin{equation}\label{eq:lin}
x\,C(x)+C(x^2)\equiv x^3\pmod m .
\end{equation}

\emph{Step 3: solving \eqref{eq:lin}.} The series $S=\sum_{j\ge0}(-1)^jx^{2^j+1}=x^2-x^3+x^5-x^9+\cdots$ satisfies
\eqref{eq:lin} exactly, because
\[
x\,S=x^3-x^4+x^6-x^{10}+x^{18}-\cdots,\qquad S(x^2)=x^4-x^6+x^{10}-x^{18}+\cdots,
\]
and everything except $x^3$ cancels. Moreover, \eqref{eq:lin} has only one solution over $\Z/m\Z$ that begins in
degree $2$. Indeed, suppose two solutions differed, and their difference $D$ began in degree $d\ge2$. Then $xD$ would
begin in degree $d+1$, while $D(x^2)$ begins in degree $2d>d+1$. So $xD+D(x^2)$ would have a nonzero coefficient in
degree $d+1$, although it must be $0$. Hence $C\equiv S\pmod m$, and $A_m=x+mC\equiv x+mS\pmod{m^2}$.
\end{proof}

\begin{proof}[Proof of Theorem~\ref{thm:binary}]
Divide both series by $x$: write $A_m=x\,F$ and $B_m=x\,G$ with
\[
F=f_0+f_1x+f_2x^2+\cdots,\qquad G=g_0+g_1x+g_2x^2+\cdots,\qquad f_0=g_0=1 .
\]
Then $f_r=a_{r+1}(m)$ and $g_r=b_{r+1}(m)$, and we must show that $m^{s_2(r)}$ divides $f_r$ and $g_r$.

\emph{Two identities.} Substitute both sides of \eqref{eq:general} into $B_m$. Since $B_m(A_m(y))=y$ for every series
$y$ without constant term, this gives $xA_m-mxA_m^2=B_m(x^2)$, which is Hanna's formula. In terms of $F$ and $G$ it reads
$x^2F-mx^3F^2=x^2G(x^2)$, that is,
\begin{equation}\label{eq:FG}
F(x)=G(x^2)+m\,x\,F(x)^2,
\qquad\text{so}\qquad
f_r=\bigl[\,r\text{ even}\,\bigr]\,g_{r/2}+m\sum_{i+j=r-1}f_if_j\quad(r\ge1),
\end{equation}
where $[\,r\text{ even}\,]$ is $1$ if $r$ is even and $0$ otherwise. The second identity is $A_m(B_m(x))=x$. Since
$A_m(y)=\sum_jf_jy^{j+1}$, it reads $\sum_{j\ge0}f_jx^{j+1}G^{j+1}=x$, that is, $\sum_jf_jx^jG^{j+1}=1$. Its
coefficient of $x^r$, for $r\ge1$, gives
\begin{equation}\label{eq:grec}
g_r=-f_r-\sum_{j=1}^{r-1}f_j\,\coeff{r-j}{G^{\,j+1}}.
\end{equation}

\emph{The induction.} We show that $m^{s_2(r)}$ divides $f_r$ and $g_r$, by induction on $r$. For $r=0$ there is
nothing to show, because $s_2(0)=0$. Let $r\ge1$, and assume the claim for all smaller indices. We first treat $f_r$,
using \eqref{eq:FG}. If $r$ is even, $g_{r/2}$ is divisible by $m^{s_2(r/2)}=m^{s_2(r)}$, by \eqref{eq:digits}. Each
term $m\,f_if_j$ with $i+j=r-1$ is divisible by $m^{1+s_2(i)+s_2(j)}$, and
$1+s_2(i)+s_2(j)\ge s_2(1+i+j)=s_2(r)$, again by \eqref{eq:digits} (note $s_2(1)=1$). So $m^{s_2(r)}$ divides $f_r$.

Now we treat $g_r$, using \eqref{eq:grec}. The term $f_r$ was just handled. Each remaining term
$f_j\,\coeff{r-j}{G^{j+1}}$ is a sum of products $f_j\,g_{i_1}g_{i_2}\cdots g_{i_{j+1}}$ with
$j+i_1+\cdots+i_{j+1}=r$ and all indices smaller than $r$. By the induction hypothesis such a product is divisible by
$m^{s_2(j)+s_2(i_1)+\cdots+s_2(i_{j+1})}$, and this exponent is at least $s_2(r)$ by \eqref{eq:digits}. So
$m^{s_2(r)}$ divides $g_r$, which completes the induction and proves~(a). For the special case, note
$s_2(2^k-1)=k$.

For (b), run the same argument over the ring $\Z[t]$, with the variable $t$ in place of $m$. Nothing changes: the
induction only uses the identities \eqref{eq:FG} and \eqref{eq:grec} and the fact that divisibility by powers of $t$
is preserved by sums and products.
\end{proof}

\begin{proof}[Proof of Corollary~\ref{cor:val}]
Since $p^{e}$ divides $m$, Theorem~\ref{thm:binary} shows that $p^{e\,s_2(n-1)}$ divides $a_n(m)$, which is the lower
bound. Now let $n\ge2$. If $n=2^j+1$, Theorem~\ref{thm:modsq} gives $a_n(m)=m\bigl((-1)^j+m\,q\bigr)$ for some integer
$q$. The second factor is $\pm1$ plus a multiple of $p$, so it is not divisible by $p$; hence $v_p(a_n(m))=v_p(m)=e$.
If $n\ge2$ is not of this form, then $n-1\ge1$ is not a power of $2$, so $s_2(n-1)\ge2$, and the lower bound gives
$v_p(a_n(m))\ge2e>e$.
\end{proof}

\begin{proof}[Proof of Corollary~\ref{cor:entry}]
Take $m=3$ in Theorems~\ref{thm:parity}, \ref{thm:modsq} and~\ref{thm:binary}. In Theorem~\ref{thm:modsq}, an index
$n=2^j+1$ has either $j=2k$ even, so that $n=4^k+1$ and the remainder is $3\cdot(+1)=3$, or $j=2k+1$ odd, so that
$n=2\cdot4^k+1$ and the remainder is $3\cdot(-1)=-3\equiv6\pmod9$.
\end{proof}

\section{Checking the results}\label{sec:checks}

A mathematical proof can contain mistakes, so the results were checked in two independent ways.

\para{Numerical check.}
We computed $A$ with the construction recorded in the entry: starting from $A=x$, repeatedly replace $A$
by $\bigl(1-\sqrt{1-12B(x^2)/x}\bigr)/6$, where $B$ is the compositional inverse of the current $A$. This reproduces all
23 terms printed in the entry, and the result satisfies the defining equation~\eqref{eq:entry} through degree $260$.
For every $n<260$, the coefficient $a(n)$ is odd, the remainders modulo $9$ are as in Corollary~\ref{cor:entry}, and
$3^{s_2(n-1)}$ divides both $a(n)$ and $b_n(3)$. For the general family, Theorems~\ref{thm:parity}, \ref{thm:modsq}
and~\ref{thm:binary}(a) and the characterization in Corollary~\ref{cor:val} hold for every $m$ from $2$ to $12$ and
every $n<70$.
The computation was done with the PARI/GP computer algebra system, using code that was written separately from the
formal proofs (file \texttt{verification/independent\_checks.gp} in the accompanying repository~\cite{repo}).
A computation can only test finitely many cases, so it is a check on the definitions and on the indexing, not a proof.

\para{Formal proof in Lean.}
Lean~\cite{lean} is a \emph{proof assistant}: a program in which mathematical statements and their proofs are written
in a precise formal language, and every step is checked mechanically down to the basic rules of logic. The theorems of Section~\ref{sec:results} are proved in Lean~4 (version 4.19.0) on top of
Mathlib~\cite{mathlib}, Lean's library of formalized mathematics. The statements are collected in the file \texttt{lean/Entries/A396843.lean} of the repository,
in the namespace \texttt{OEIS.A396843}. There the left side of \eqref{eq:general} is \texttt{Series.Psi m A}, which stands for
\texttt{compose A (X*A - C m * X * A\^{}2)}; the coefficients $a_n(m)$ and $b_n(m)$ are \texttt{a m n} and
\texttt{b m n}. Sixteen principal statements about this family are collected in
\texttt{lean/Hanna/Statements.lean}.
The table lists where each result of this paper is proved.
\begin{center}\small\renewcommand{\arraystretch}{1.25}
\begin{tabular}{>{\raggedright\arraybackslash}p{0.52\linewidth}>{\raggedright\arraybackslash}p{0.4\linewidth}}
\toprule
Statement in this paper & Lean theorem(s)\\
\midrule
Lemma~\ref{lem:exist} (every commutative ring, every $m$) & \texttt{existence\_\allowbreak{}uniqueness} \\
Lemma~\ref{lem:inverse} (compositional inverse) & \texttt{compositional\_\allowbreak{}inverse} \\
$a_n(m)$ and $b_n(m)$ are values of integer polynomials in $m$ & \texttt{coefficient\_\allowbreak{}polynomial\_\allowbreak{}specialization} \\
Theorem~\ref{thm:parity} & \texttt{parity\_\allowbreak{}odd}, \texttt{parity\_\allowbreak{}even} \\
Theorem~\ref{thm:modsq} & \texttt{mod\_\allowbreak{}parameter\_\allowbreak{}squared} \\
Theorem~\ref{thm:binary}, parts (a) and (b) & \texttt{binary\_\allowbreak{}divisibility\_\allowbreak{}integer}, \texttt{inverse\_\allowbreak{}binary\_\allowbreak{}divisibility\_\allowbreak{}integer}, \texttt{binary\_\allowbreak{}divisibility\_\allowbreak{}polynomial}, \texttt{inverse\_\allowbreak{}binary\_\allowbreak{}divisibility\_\allowbreak{}polynomial}, \texttt{power\_\allowbreak{}of\_\allowbreak{}two\_\allowbreak{}index} \\
Corollary~\ref{cor:val} & \texttt{prime\_\allowbreak{}divisibility\_\allowbreak{}lower\_\allowbreak{}bound}, \texttt{prime\_\allowbreak{}exponent\_\allowbreak{}exists\_\allowbreak{}unique}, \texttt{minimum\_\allowbreak{}prime\_\allowbreak{}exponent\_\allowbreak{}iff} \\
Corollary~\ref{cor:entry}, and its analogue for $m=4$ & \texttt{a396843\_\allowbreak{}parity}, \texttt{residues\_\allowbreak{}parameter\_\allowbreak{}three}, \texttt{residues\_\allowbreak{}parameter\_\allowbreak{}four} \\
\bottomrule
\end{tabular}
\end{center}
The Lean project was rebuilt from source, and the axioms on which each of these theorems depends were listed with
Lean's \texttt{\#print axioms} command. Every one depends only on the three standard axioms of Lean's logic
(\texttt{propext}, \texttt{Classical.choice}, \texttt{Quot.sound}); no step is left unproved. A formal proof certifies
exactly the statement that Lean checks. Whether that statement says what the OEIS entry means is a separate question,
so the Lean definitions were also compared with the definition in the OEIS entry
(Section~\ref{sec:sequence}).

\section{Discussion}\label{sec:discussion}

\para{Scope.} The results settle the first conjecture of the entry, and replace the second by a complete and proved
description of the remainders modulo $9$. They say nothing about the size of the coefficients, or about the radius of
convergence that the entry also discusses.

\para{Higher powers of $m$.} Theorem~\ref{thm:binary} also gives information modulo every power of $m$: modulo
$m^k$, the coefficient $a_n(m)$ is $0$ unless $n-1$ has fewer than $k$ ones in its binary expansion. So modulo $m^k$
the nonzero remainders can occur only at a very sparse set of indices. For $k=2$, Theorem~\ref{thm:modsq} describes these
remainders completely. For $k=3$ the candidates are $n=1$, $n=2^i+1$ and $n=2^i+2^j+1$ with $i>j$. The remainders modulo
$m^3$ are not determined here.

\para{Other values of $m$.} For $m=1$ (entry A265940) the results say only that every coefficient is odd. For $m=2$
(entry A271958) they say, for example, that $a_n(2)\equiv2\pmod 4$ when $n=2^j+1$ and $a_n(2)\equiv0\pmod4$ for the
other $n\ge2$. For $m=5$ (entry A396845) the remainders modulo $25$ are $5$ and $20$ at $n=4^k+1$ and $n=2\cdot4^k+1$,
respectively, and $0$ for the other $n\ge2$.

\section*{About this paper}\label{sec:about}
\para{How the work was done.}
This paper is one of 29 papers on conjectures from the OEIS that came out of an AI research effort directed by
Dev~D.~Goyal~\cite{repo}. The roles were as follows.
\begin{itemize}
\item \emph{Dev~D.~Goyal} started the project, chose which problems to pursue, asked for proofs, formal verification and
revisions, and decided what to publish. Dev~D.~Goyal did not supply proof strategies or intermediate arguments.
\item \emph{ChatGPT} (OpenAI), using the GPT-6 model with Pro reasoning, worked autonomously from those prompts. It
searched the OEIS for suitable conjectures, found the proofs, wrote the Lean~4 formalization, and wrote a first draft
of this paper.
\item \emph{Claude Opus~5.5}, used with Claude Code, checked the work independently. It rebuilt the Lean project from source and confirmed that the theorems listed in Section~\ref{sec:checks} rely only on
Lean's standard axioms; compared the Lean definitions with the OEIS definition; recomputed the sequence with separately
written PARI/GP code, compared it with the terms published in the OEIS, and tested every claim of the paper on many
terms; and read every step of the proofs. It then wrote the present version of the paper, expanding the
explanations and, in places, writing out an argument in more detail or by a more elementary route. The theorems of
Section~\ref{sec:results} and the proof strategies are ChatGPT's. The Discussion section, including the short arguments
in it, was written by Claude Opus~5.5; its observations were checked by computer and are not part of the Lean
formalization.
\end{itemize}
Dev~D.~Goyal has verified the results and is responsible for this paper.

\para{Review status.} This paper has not been peer reviewed, and it has not been submitted to a journal or to the OEIS.

\para{Priority.} An OEIS entry that labels a statement as a conjecture shows only that its contributor did not know
a proof; it does not show that no proof existed. During the research only limited searches for earlier work were
made. No claim of priority is made.

\para{Corrections.} Please report any error in the issue tracker of the repository~\cite{repo}, where the Lean sources and the numerical checks are also
available.

\section*{Acknowledgements}
The family of sequences and the conjectures are due to Paul~D.~Hanna~\cite{oeis1,oeis2,oeis3,oeis4,oeis5}. The identity
$xA-mxA^2=B(x^2)$, used in the proof of Theorem~\ref{thm:binary}, is recorded in his entries. Danial~Asaria found the counterexample $a(5)=975$ to the conjecture modulo~$9$~\cite{oeis3}.

\begin{thebibliography}{99}\small
\raggedright
\bibitem{oeis1} P.~D.~Hanna, entry A265940 (G.f.\ $A(x)$ satisfies $A(xA(x)-xA(x)^2)=x^2$), \emph{The On-Line Encyclopedia of Integer Sequences}, 2015. \url{https://oeis.org/A265940}. Text as of October 2026.
\bibitem{oeis2} P.~D.~Hanna, entry A271958 (G.f.\ $A(x)$ satisfies $A(xA(x)-2xA(x)^2)=x^2$), \emph{The On-Line Encyclopedia of Integer Sequences}, 2016. \url{https://oeis.org/A271958}. Text as of October 2026.
\bibitem{oeis3} P.~D.~Hanna, entry A396843 (G.f.\ $A(x)$ satisfies $A(xA(x)-3xA(x)^2)=x^2$), \emph{The On-Line Encyclopedia of Integer Sequences}, July 2026. \url{https://oeis.org/A396843}. Text as of October 2026.
\bibitem{oeis4} P.~D.~Hanna, entry A396844 (G.f.\ $A(x)$ satisfies $A(xA(x)-4xA(x)^2)=x^2$), \emph{The On-Line Encyclopedia of Integer Sequences}, July 2026. \url{https://oeis.org/A396844}. Text as of October 2026.
\bibitem{oeis5} P.~D.~Hanna, entry A396845 (G.f.\ $A(x)$ satisfies $A(xA(x)-5xA(x)^2)=x^2$), \emph{The On-Line Encyclopedia of Integer Sequences}, July 2026. \url{https://oeis.org/A396845}. Text as of October 2026.
\bibitem{concrete} R.~L.~Graham, D.~E.~Knuth and O.~Patashnik, \emph{Concrete Mathematics}, second edition, Addison-Wesley, 1994. (Section~4.4 proves Legendre's formula.)
\bibitem{oeisdb} N.~J.~A.~Sloane and The OEIS Foundation Inc., \emph{The On-Line Encyclopedia of Integer Sequences}. \url{https://oeis.org}.
\bibitem{stanley} R.~P.~Stanley, \emph{Enumerative Combinatorics}, Volume~2, Cambridge University Press, 1999. (Chapter~5 treats formal power series, composition and compositional inverses.)
\bibitem{niven} I.~Niven, Formal power series, \emph{American Mathematical Monthly} 76 (1969), 871--889.
\bibitem{lean} L.~de~Moura and S.~Ullrich, The Lean~4 theorem prover and programming language, in \emph{Automated Deduction -- CADE 28}, Lecture Notes in Computer Science 12699, Springer, 2021, 625--635.
\bibitem{mathlib} The mathlib Community, The Lean mathematical library, in \emph{Proceedings of the 9th ACM SIGPLAN International Conference on Certified Programs and Proofs (CPP 2020)}, 367--381. Version used: commit \texttt{c44e0c8ee63c}.
\bibitem{repo} D.~D.~Goyal, \emph{Proofs and corrections of OEIS conjectures}, repository, 2026. \url{https://github.com/dvgyl/oeis-a396843-proof}.
\end{thebibliography}
\end{document}
